\documentclass[12pt, a4paper, twoside]{book}

% ─── Packages ───────────────────────────────────────────────────────────────
\usepackage[a4paper, inner=3.0cm, outer=2.5cm, top=2.8cm, bottom=2.8cm, includeheadfoot, headheight=24pt]{geometry}
\usepackage[T1]{fontenc}
\usepackage[utf8]{inputenc}
\usepackage{amsmath, amsthm}
\usepackage[charter, cal=cmcal]{mathdesign}
\usepackage[scaled=0.92]{helvet}
\usepackage{mathtools}
\usepackage{physics}
\usepackage{graphicx}
\usepackage{xcolor}
\usepackage{tikz}
\usetikzlibrary{arrows.meta, decorations.pathmorphing, shapes.geometric,
                positioning, calc, patterns, backgrounds, decorations.markings}
\usepackage{pgfplots}
\pgfplotsset{compat=1.18}
\usepackage{tcolorbox}
\tcbuselibrary{skins, breakable, theorems}
\usepackage{enumitem}
\usepackage{booktabs}
\usepackage{array}
\usepackage{multirow}
\usepackage{fancyhdr}
\usepackage{titlesec}
\usepackage{hyperref}
\usepackage{wrapfig}
\usepackage{float}
\usepackage{caption}
\usepackage{subcaption}
\usepackage{epigraph}
\usepackage{lettrine}
\usepackage{mdframed}
\usepackage{setspace}
\usepackage{parskip}
\usepackage{siunitx}

% ─── Colours ────────────────────────────────────────────────────────────────
\definecolor{deep}{HTML}{0F172A}      % Deep Slate/Charcoal for body text/titles
\definecolor{accent}{HTML}{E11D48}    % Crimson/Coral for problem counters
\definecolor{teal}{HTML}{0D9488}      % Electric Teal for rules and highlights
\definecolor{gold}{HTML}{D97706}      % Gold for puzzle borders
\definecolor{midblue}{HTML}{334155}   % Muted Blue-Gray for section headings
\definecolor{purple}{HTML}{6D28D9}    % Purple for reflection borders
\definecolor{lightbg}{HTML}{F8FAFC}   % Very light gray-blue for exploration bg
\definecolor{warmbg}{HTML}{FEF3C7}    % Warm gold-yellow for puzzle bg
\definecolor{greenbg}{HTML}{F0FDF4}   % Light green for thought/keyidea bg

% ─── Page geometry fix ──────────────────────────────────────────────────────
\setlength{\headheight}{24pt}

% ─── Hyperref ───────────────────────────────────────────────────────────────
\hypersetup{
  colorlinks=true, linkcolor=midblue, urlcolor=accent, citecolor=teal,
  pdftitle={Discovery Physics: Journey 2 — Motion and Its Causes},
  pdfauthor={Discovery Institute},
}

% ─── Page style ─────────────────────────────────────────────────────────────
\pagestyle{fancy}
\fancyhf{}
\fancyhead[LE]{\small\textcolor{deep}{\textit{Discovery Physics}}}
\fancyhead[RO]{\small\textcolor{deep}{\textit{\leftmark}}}
\fancyfoot[C]{\small\textcolor{deep}{\thepage}}
\renewcommand{\headrulewidth}{0.4pt}
\renewcommand{\headrule}{\hbox to\headwidth{\color{teal}\leaders\hrule height \headrulewidth\hfill}}

% ─── Chapter / Section titles ───────────────────────────────────────────────
\titleformat{\chapter}[display]
  {\normalfont\huge\bfseries\color{deep}}
  {\textcolor{accent}{\chaptertitlename\ \thechapter}}{12pt}{\Huge}
  [\vspace{4pt}\textcolor{teal}{\titlerule[1.5pt]}]

\titleformat{\section}
  {\normalfont\Large\bfseries\color{midblue}}
  {\textcolor{accent}{\thesection}}{1em}{}

\titleformat{\subsection}
  {\normalfont\large\bfseries\color{deep}}
  {\textcolor{teal}{\thesubsection}}{1em}{}

% ─── Custom boxes ───────────────────────────────────────────────────────────
\newtcolorbox{puzzlebox}[1][]{
  enhanced, breakable,
  colback=warmbg, colframe=gold,
  fonttitle=\bfseries\large\color{deep}, title={The Puzzle},
  attach boxed title to top left={yshift=-2mm, xshift=4mm},
  boxed title style={colback=gold, colframe=gold}, #1}

\newtcolorbox{explobox}[1][]{
  enhanced, breakable,
  colback=lightbg, colframe=midblue,
  fonttitle=\bfseries\color{white}, coltitle=white, title={Exploration},
  attach boxed title to top left={yshift=-2mm, xshift=4mm},
  boxed title style={colback=midblue, colframe=midblue}, #1}

\newtcolorbox{discoverybox}[1][]{
  enhanced, breakable,
  colback=white, colframe=accent,
  fonttitle=\bfseries\large\color{white}, coltitle=white, title={Discovery},
  attach boxed title to top left={yshift=-2mm, xshift=4mm},
  boxed title style={colback=accent, colframe=accent}, drop shadow, #1}

\newtcolorbox{namebox}[1][]{
  enhanced,
  colback=deep!5, colframe=deep,
  fonttitle=\bfseries\color{white}, coltitle=white, title={What This Is Called},
  attach boxed title to top left={yshift=-2mm, xshift=4mm},
  boxed title style={colback=deep, colframe=deep}, #1}

\newtcolorbox{reflectionbox}[1][]{
  enhanced, breakable,
  colback=purple!5, colframe=purple,
  fonttitle=\bfseries\color{white}, coltitle=white, title={Reflection}, #1}

\newtcolorbox{historybox}[1][]{
  enhanced, breakable,
  colback=gold!8, colframe=gold!70!brown,
  fonttitle=\bfseries\color{deep}, title={Historical Note}, #1}

% ─── Problem counter ────────────────────────────────────────────────────────
\newcounter{problem}[chapter]
\newcommand{\problem}[1][]{%
  \refstepcounter{problem}%
  \medskip\noindent
  \textcolor{accent}{\textbf{Problem \theproblem.}}\ %
  \ifx&#1&\else\textit{(#1)}\ \fi}

\newcommand{\keyidea}[1]{%
  \begin{tcolorbox}[enhanced, colback=teal!8, colframe=teal,
    left=4mm, fontupper=\bfseries\color{deep}]#1\end{tcolorbox}}

% ─── Theorem environments ───────────────────────────────────────────────────
\newtheoremstyle{disc}{8pt}{8pt}{\itshape}{}{\bfseries\color{midblue}}{.}{.5em}{}
\theoremstyle{disc}
\newtheorem{principle}{Principle}[chapter]
\newtheorem{question}{Question}[chapter]

% ─── Spacing ────────────────────────────────────────────────────────────────
\setlength{\epigraphwidth}{0.65\textwidth}
\renewcommand{\epigraphflush}{center}
\setlength{\parskip}{6pt}
\setlength{\parindent}{0pt}
\onehalfspacing

% ════════════════════════════════════════════════════════════════════════════
\begin{document}


\setcounter{chapter}{1}
\setcounter{section}{5}

\section{Module 2.6 — The Moon Is Falling}
\label{sec:moon_falling}

\epigraph{
  \itshape I began to think of gravity extending to the orb of the Moon.
}{--- Isaac Newton, recalling 1666}

\subsection{The Central Question}

\begin{tcolorbox}[colback=deep!5, colframe=deep, fontupper=\large\bfseries\color{deep}]
Is the Moon falling toward Earth? And is the force that holds it in orbit the same force that pulls an apple from a tree?
\end{tcolorbox}

\begin{figure}[H]
  \centering
  \includegraphics[width=0.9\textwidth]{figures/fig_j2_moon}
  \caption{Left: the Moon's orbit as continuous free fall. The velocity is tangential; gravity pulls it centripetally inward. Without gravity, it would fly off in a straight line. Right: the numerical comparison showing that the Moon's centripetal acceleration agrees exactly with the inverse-square prediction from surface gravity.}
  \label{fig:moon}
\end{figure}

\subsection{Newton's Insight}

The story (probably embellished, but instructive) is that Newton saw an apple fall and asked: how far does gravity reach? Could it reach the Moon?

At the Moon's distance (about 60 Earth radii), if gravity follows an inverse-square law, the gravitational acceleration should be:
\begin{equation}
  g_{\text{Moon}} = \frac{g_{\text{surface}}}{60^2} = \frac{\SI{9.81}{m/s^2}}{3600} \approx \SI{0.00272}{m/s^2}
\end{equation}

Now we can check: does the Moon's observed orbital motion require exactly this centripetal acceleration?

\begin{explobox}
\problem In Module 2.3, Problem 2, you computed the Moon's centripetal acceleration from its orbital data. What did you find? Compare it to the value predicted above from the inverse-square law.

\problem The result should match to within the precision of the input data. Newton reportedly found that his calculation ``answered pretty nearly.'' Why was he cautious about claiming an exact match? What sources of uncertainty existed?

\problem Now work backward. From the fact that the Moon's acceleration is $\SI{0.00272}{m/s^2}$ at 60 Earth radii, and the surface acceleration is $\SI{9.81}{m/s^2}$, what can you deduce about how gravity varies with distance?

\problem The Sun's gravitational pull on Earth is about 180 times stronger than the Moon's. Yet the Moon dominates the tides on Earth. (Ocean tides are caused by the \textit{difference} in gravitational force across Earth's diameter.) Explain qualitatively why the Moon dominates despite the Sun being so much more massive.
\end{explobox}

\subsection{The Universal Law of Gravitation}

Newton's great generalisation: the gravitational force between any two masses $M$ and $m$ separated by distance $r$ is:

\begin{equation}
  \boxed{F = \frac{GMm}{r^2}}
\end{equation}

where $G = \SI{6.674e-11}{N.m^2/kg^2}$ is the gravitational constant, measured by Henry Cavendish in 1798---more than a century after Newton formulated the law.

\begin{explobox}
\problem Calculate the gravitational force between:
\begin{enumerate}[label=(\alph*)]
  \item Two $\SI{1}{kg}$ masses separated by $\SI{1}{m}$ (compare to their weight)
  \item Earth ($\SI{6.0e24}{kg}$) and the Moon ($\SI{7.3e22}{kg}$) at their actual distance
  \item The Sun ($\SI{2.0e30}{kg}$) and Earth at $\SI{1.5e11}{m}$
\end{enumerate}

\problem At what height above Earth's surface is the gravitational acceleration exactly half of its surface value?

\problem The gravitational acceleration on the surface of the Moon is $\SI{1.62}{m/s^2}$. The Moon's radius is $\SI{1.74e6}{m}$. Calculate the Moon's mass.

\problem \textit{Kepler from Newton.} Show that Newton's law of gravitation, combined with $F = ma$ for circular orbits, gives Kepler's third law: $T^2 \propto r^3$.
\end{explobox}

\begin{discoverybox}
\textbf{The unification.} The same formula---$F = GMm/r^2$---governs:
\begin{itemize}
  \item An apple falling 1 metre to the ground
  \item The Moon orbiting Earth at $\SI{384000}{km}$
  \item Earth orbiting the Sun at $\SI{150}{Gm}$
  \item A distant galaxy bending light around its mass
\end{itemize}

This was Newton's true achievement. Not the formula itself---Hooke, Halley, and others had guessed the inverse-square form. The achievement was the \textit{proof} that this single formula, combined with $F = ma$, explains all of Kepler's laws and all terrestrial mechanics with the same mathematics.

Heaven and Earth obey the same physics. This was not obvious before 1687. It has been confirmed to extraordinary precision since.
\end{discoverybox}

% ────────────────────────────────────────────────────────────────────────────

\end{document}
